%% kapa: exam=1 round=37 year=2026 period=1 form=A subject=economics subject_ko=경제학원론 seq=20 number=60
%% source: Q-Net 감정평가사 기출문제 (cst00309) artlSeq=5258924 file=raw/1008/5258924/2026년 제37회 감정평가사 1차 1교시.pdf
%% conversion: kordoc 4.24.51 -> tools/gen1.py (auto); review=proofread-sonnet; flags=none
\begin{question}{60}
정보의 비대칭성에 관한 설명으로 옳은 것은?
\begin{choices}
  \item 사고가 발생할 가능성이 높은 사람일수록 보험에 가입할 가능성이 크다는 것은 도덕적해이(moral hazard)에 해당한다.
  \item 화재보험 가입자가 화재예방 노력을 게을리할 가능성이 크다는 것은 역선택(adverse selection)에 해당한다.
  \item 선별(screening)은 정보를 보유하지 못한 측이 도덕적해이 문제를 해결하기 위해 사용할 수 있는 방법이다.
  \item 항공사가 서로 다른 유형의 소비자에게 각각 다른 요금을 부과하는 행위는 신호발송(signaling)에 해당한다.
  \item 작업감독제는 도덕적해이 방지를 위한 수단 중 하나이다.
\end{choices}
\end{question}
\answer{5}
